\documentclass[11pt]{article}
\usepackage[utf8]{inputenc}
\usepackage[T1]{fontenc}
\usepackage{lmodern}
\usepackage[a4paper,margin=2.4cm]{geometry}
\usepackage{xcolor}
\usepackage{framed}
\usepackage{amsmath}
\usepackage{enumitem}
\usepackage{hyperref}
\hypersetup{colorlinks=true,linkcolor=blue,urlcolor=blue}

\definecolor{commentbg}{RGB}{238,242,248}
\definecolor{commentrule}{RGB}{90,120,170}

\newenvironment{comment}{%
  \def\FrameCommand{{\color{commentrule}\vrule width 3pt}\hspace{0pt}%
    \fboxsep=8pt\colorbox{commentbg}}%
  \MakeFramed{\advance\hsize-\width\FrameRestore}%
  \itshape\small}{\endMakeFramed}

\newcommand{\reviewerpoint}[1]{\medskip\noindent\textbf{Comment #1.}\;}
\newcommand{\response}{\par\smallskip\noindent\textbf{Response:}\;}

\title{\vspace{-1.5em}Response to Reviewer 1}
\author{}
\date{}

\begin{document}
\maketitle
\vspace{-3em}

\noindent We thank the reviewer for a careful and constructive reading of the
manuscript. The comments on novelty framing, the oracle-efficiency metric, the
fairness of baseline comparisons, the strength of the causal claims, and the
clarity of the figures have substantially improved the paper. The manuscript has
been revised throughout; all changes are summarised below, with pointers to the
relevant sections, tables, and figures of the revised manuscript. Page numbers
refer to the revised compiled PDF.

\section*{Summary of the main changes}
\begin{itemize}[leftmargin=1.4em,itemsep=2pt]
  \item The contribution is now framed explicitly as an intervention-selection
        layer on top of an existing attribution backend, not a new attribution
        method (Introduction, p.~2; Section~``Comparison with Published Circuit
        Discovery Methods'', Table~11, p.~31).
  \item The $>100\%$ ``oracle efficiency'' has been replaced by two clearly
        separated metrics: a bounded oracle efficiency ($\leq 100\%$) for
        ablation-only selectors, and a Relative Cumulative KL (RCK) for
        steering-enabled selectors (Section~``Metrics'', p.~16--17).
  \item A dose--response steering sweep and a random-feature control were added,
        and the causal claims were weakened accordingly (Section~``Feature
        Steering'', p.~13; Results, p.~24).
  \item Figures~1 and~2 were redrawn as clean raster diagrams; the remaining
        figures were given shared, non-overlapping legends.
  \item The oracle rows of the KL tables now report the actual oracle KL values.
  \item Action-matched (steering-enabled) variants of random, greedy, and bandit
        were added, plus a random-action control, and the Llama multi-step
        shortfall is now treated as a limitation.
  \item A plain UCB baseline, a cross-task head-to-head efficiency table
        (Table~9, p.~29), a hyperparameter-sensitivity table (Table~8, p.~29),
        and an online belief-convergence figure (Figure~9, p.~24) were added to
        strengthen the empirical comparison.
\end{itemize}

\section*{Point-by-point responses}

\begin{comment}
\reviewerpoint{1} The novelty of the proposed framework should be clarified.
It is not fully clear whether the main contribution is a new circuit discovery
method or an active intervention-selection strategy built on existing attribution
graph tools.
\end{comment}
\response The contribution is the latter, and the manuscript now states this
plainly. The revised Introduction (p.~2) describes the work as ``an
intervention-selection layer that sits on top of an existing attribution
backend'', and the comparison section (Table~11, p.~31) explicitly lists
\texttt{circuit-tracer} as the attribution backend and
positions the agent as the selection policy operating over its candidate set.
The novelty claimed is the active-inference POMDP formulation of which
feature and which intervention type to test next under a fixed budget, not the
construction of the attribution graph itself.

\begin{comment}
\reviewerpoint{2} The ``oracle efficiency'' metric needs a clearer definition.
Reporting values above 100\%, such as 1255\%, is confusing if the oracle is an
upper bound. This should be mathematically explained or renamed.
\end{comment}
\response The metric has been split and renamed. For ablation-only selectors,
every method and the oracle draw from the same quantity (ablation KL), so the
ratio is a bounded efficiency $\eta\in[0,100]\%$
(Equation for $\eta$, Section ``Metrics'', p.~16). For the full multi-action
agent, which may steer, the cumulative KL is compared with the ablation oracle
through a separate quantity, the Relative Cumulative KL
(RCK, Equation for RCK, p.~17), explicitly labelled as a relative amplification
factor and never as an efficiency. The 1255\% figure now appears only as an RCK
value in Table~3 (p.~22) and is described as steering-driven amplification rather
than super-oracle discovery.

\begin{comment}
\reviewerpoint{3} The causal claims should be stated more cautiously. Feature
steering with 10$\times$ activation scaling may cause out-of-distribution
effects, so additional controls such as smaller scaling factors and
random-feature comparisons would strengthen the claim.
\end{comment}
\response Both controls were added. The steering protocol now sweeps
$m\in\{0,1.5,2,3,5,10\}$ and compares each circuit-selected feature against a
matched random-feature control steered at the same multipliers
(Section ``Feature Steering'', p.~13). The Results section (p.~24) reports that
the selectivity advantage of circuit features over the control is not
statistically significant at large multipliers (one-sided Fisher exact test), so
the steering result is now framed as a controllability demonstration rather than
proof of privileged selectivity. A complementary top-$k$ token analysis shows
that large multipliers amplify coherent concept tokens rather than only
degrading the model, which is reported alongside the dose--response curve.

\begin{comment}
\reviewerpoint{4} Figure~1 should be visually revised. Some arrows/lines appear
to cross through each other, which makes the workflow difficult to follow.
\end{comment}
\response Figures~1 and~2 (system architecture and the EFE pipeline) were
re-authored as clean raster diagrams (\texttt{fig1.png}, \texttt{fig2.png}) with
non-crossing routing and clear spacing, replacing the previous TikZ figures
(Methodology, p.~8 and p.~9).

\begin{comment}
\reviewerpoint{5} In Tables~2 and~4, the oracle row reports only 100\% efficiency
but not the corresponding mean or cumulative KL value. The authors should report
the actual oracle KL scores to make the normalization transparent.
\end{comment}
\response The oracle KL values are now reported. Table~2 (IOI, p.~19) lists the
oracle's per-prompt mean KL (0.000862 for Gemma, 0.009755 for Llama), and the
surrounding text gives the cumulative-KL denominators (0.0172 and 0.1951). The
multi-step table (Table~5, p.~26) likewise reports the oracle cumulative KL (0.0101 Gemma,
0.2447 Llama). The normalisation underlying every efficiency value is therefore
explicit.

\begin{comment}
\reviewerpoint{6} The POMDP agent performs worse than the bandit in the Llama
multi-step test (33.6\% vs.\ 76.4\%). This negative result is significant and
should be treated as a limitation, not an architectural explanation.
\end{comment}
\response The Llama multi-step shortfall is now reported as a limitation.
Under the fair ablation-only protocol the multi-step results appear in Table~5
(p.~26), and the Discussion (``Cross-Model Generality'' and ``Limitations'',
p.~31) states that the agent underperforms on Llama multi-step and that a
finer-layer-bin variant was tried but did not recover the gap (Table~5). The
result is presented as an open limitation tied to the coarse layer-role
discretisation rather than as a feature of the architecture.

\begin{comment}
\reviewerpoint{7} Only an ablation version of the POMDP is included; action-matched
variants are missing for random, greedy, and bandit. Does the POMDP's success stem
from EFE-guided selection or from the steering action producing higher KL?
Action-matched versions of the baselines would make the comparison more balanced.
\end{comment}
\response Action-matched baselines were added. Greedy+steer, Bandit+steer, and
Random+steer use the identical multi-action space as the full agent, and a
Random-action control randomises the intervention type
(Section ``Baselines'', p.~15). Reported on the common RCK scale (Table~3, p.~22),
these show that on Gemma every steering-enabled strategy, including random+steer
(888\%), produces RCK well above 100\%, so the amplification is attributable to
the steering action itself rather than to which feature is selected. This
decomposition is exactly what isolates EFE-guided selection from the steering
effect, and is discussed in Section ``Multi-Action Amplification'' (p.~22). A
plain UCB baseline (textbook UCB1 over layers, using observed KL as reward) was
also added to separate the contribution of the heuristic bandit's hand-built
importance and locality priors from upper-confidence exploration alone, and a
cross-task head-to-head comparison of all ablation-only selectors is reported in
Table~9 (p.~29).

\begin{comment}
\reviewerpoint{8} In Gemma IOI the agent selects feature steering in 94 of 100
steps after the initial stage, suggesting the high efficiency stems from repeated
orientation rather than balanced exploration. The authors should discuss whether
the framework is discovering circuit structure or exploiting steering-induced KL
increases.
\end{comment}
\response This distinction is now central to the analysis. The bounded oracle
efficiency (the discovery-quality claim) is computed only from the ablation-only
agent against the ablation oracle, so it cannot be inflated by steering. The
steering-driven amplification is reported separately as RCK and explicitly
described as exploitation of the steering action, not structure discovery
(Section ``Multi-Action Amplification'', p.~22; Discussion ``Structure Discovery
versus Steering-KL Exploitation'', p.~32). The action-distribution figure
(Figure~10, p.~25) shows the 94/100 steering pattern and is used to explain the
RCK contrast between the two models.

\begin{comment}
\reviewerpoint{9} Appendix~B gives fixed transition probabilities (ablation 50\%
stay / 25\% down / 25\% up; patching 70\% stay; steering 90\% stay) that are not
learned from data. The authors should justify these values.
\end{comment}
\response The transition probabilities are design choices encoding how strongly
each action is expected to move the latent importance state, and the appendix now
states this rationale: ablation can reveal or remove importance (broad, near-
symmetric transitions), patching is more conservative, and steering is least
likely to change the underlying importance estimate (concentrated diagonal). The
full $B$ matrices for every action are now given explicitly in the appendix
(``Complete Generative Model Specification'', p.~37), and a $\pm 20\%$ sensitivity
sweep over the discretisation thresholds (Table~7, p.~28) together with a
hyperparameter sweep over the agent's core parameters (Table~8, p.~29) shows the
qualitative ordering is preserved, so the conclusions are not artefacts of these
fixed values. The online Dirichlet learning is also shown to converge: Figure~9
(p.~24) plots the mean per-step L1 drift of the KL-likelihood matrix decreasing
across IOI prompts, confirming that the learned observation model stabilises
within the intervention budget.

\section*{English language and style}
The manuscript has been copy-edited throughout for concision and a consistent
third-person scientific register, and overly long sentences flagged in review
were split.

\end{document}
